\section*{الفصل \ref{ch:b2:field-lab-methods-2} --- طرائق الحقل والمختبر II}
% numbers from the Book 4 second-writer scratch calculator calc/ch31.py

\begin{solution}{exo:b2:field-lab-methods-2:1}
(أ) (1) \omterm{def:b2:field-lab-methods-2:textures}{نسيج أوفيتي}؛ (2) \omterm{def:b2:field-lab-methods-2:textures}{نسيج غرانوبلاستي}؛ (3) بويكيلوبلاست، أي بورفيروبلاست مليء
بالمحتبَسات؛ (4) هالة، أي حافّة تفاعلية حول غارنت كان يتفكّك. (ب) بحسب
\eqref{eq:b2:field-lab-methods-2:lami}، $\gamma_1/\gamma_2 = \sin\ang{104}/\sin\ang{128} =
0.9703/0.7880 = 1.23$: فذلك الحدّ يشدّ بقوّة تزيد \qty{23}{\percent} على كلّ من الحدّين الآخرين.
\end{solution}

\begin{solution}{exo:b2:field-lab-methods-2:2}
(أ) $Q + A + P = 80$: \mbox{\textdir TLT Q $= 100 \times 25/80 = 31.3$}، و\mbox{\textdir TLT A $= 18.8$}، و\mbox{\textdir TLT P $= 50.0$}. (ب) \mbox{\textdir TLT P/(A+P) $=
40/55 = \qty{72.7}{\percent}$}، بين 65 و90 مع Q بين 20 و60: إنّه غرانوديوريت.
(ج) $M' = 12 + 8 = \qty{20}{\percent}$، دون 35: فاتح اللون.
\end{solution}

\begin{solution}{exo:b2:field-lab-methods-2:3}
(أ) $\sigma = \sqrt{0.12 \times 0.88/800} = 0.0115$، أي $12.0 \pm 1.1$ نقطة مئوية، بخطأ
نسبي قدره $0.0115/0.12 = \qty{9.6}{\percent}$. (ب) أربعة أضعاف النقاط تنصّفه: $0.0057$.
\end{solution}

\begin{solution}{exo:b2:field-lab-methods-2:4}
(أ) الفقد \qty{0.0288}{g}، و$0.0288/0.8000 = \qty{3.60}{\percent}$.
(ب) $0.7712/8.7712 = 0.0879$.
(ج) المسحوق المحروق: $5.20/0.08792 = \qty{59.1}{\percent}$؛ والمسحوق الأصلي:
$59.14 \times (1 - 0.0360) = \qty{57.0}{\percent}$.
\end{solution}

\begin{solution}{exo:b2:field-lab-methods-2:5}
(أ) $\sigma = 0.10 \times 0.02 = 0.002$ و$n = 0.02 \times 0.98/0.002^2 = 4900$ نقطة.
(ب) $\sigma = 0.03$ و$n = 0.30 \times 0.70/0.03^2 = 233.3$، أي 234 نقطة. (ج) لخطأ نسبي
ثابت، يكبر $n = (1-p)/(p\,\varepsilon^2)$ بمعدّل $1/p$: فالمعدن الثانوي يتطلّب 21 ضعفًا من
النقاط، أكثر ممّا توفّره الشريحة عادةً.
\end{solution}

\begin{solution}{exo:b2:field-lab-methods-2:6}
(أ) RGM-1: $72.89/73.45 = \qty{99.2}{\percent}$؛ وBCR-2: $54.17/54.1 = \qty{100.1}{\percent}$.
(ب) RGM-1: $0.19/\sqrt{20} = 0.042$، و$-0.56/0.042 = -13.2$؛ وBCR-2: $0.13/\sqrt{20} = 0.029$،
و$0.07/0.029 = 2.4$. (ج) نتيجة RGM-1 منحازة: فمتوسّطها يبتعد عن القيمة الموصى بها 13 مرّة قدر
انحرافها المعياري، وهذا أكثر بكثير ممّا تسمح به الصدفة (إن كانت القيمة الموصى بها صحيحة). ونتيجة BCR-2 عند
حدّ ما تعطيه الصدفة. والطريقة ليست أقلّ دقّة في RGM-1: فانحرافها المعياري النسبي،
$0.19/72.89 = \qty{0.26}{\percent}$، قريب من انحراف BCR-2، \qty{0.24}{\percent}. فالمشكلة في
الصحّة، لا في الدقّة.
\end{solution}

\begin{solution}{exo:b2:field-lab-methods-2:7}
(أ) $M = 39.098 + 3 \times 24.305 + 26.982 + 3 \times 28.085 + 12 \times 15.999 + 2 \times 1.008 =
\qty{417.25}{g.mol^{-1}}$؛ ومجموعتا OH تكوّنان جزيء \ce{H2O} واحدًا ($18.015$)، أي \qty{4.32}{\percent}.
(ب) \ce{K2O} $47.10/417.25 = \qty{11.3}{\percent}$؛ و\ce{MgO} $3 \times 40.304/417.25 =
\qty{29.0}{\percent}$؛ و\ce{Al2O3} $50.98/417.25 = \qty{12.2}{\percent}$؛ و\ce{SiO2} $3 \times
60.083/417.25 = \qty{43.2}{\percent}$؛ والمجموع \qty{95.7}{\percent}. (ج) لا يكشف \omterm{def:b2:field-lab-methods-2:epma}{المسبار الإلكتروني الدقيق}
الهيدروجين: فتحليل الميكا ينقصه ماؤها، نحو \qty{4}{\percent}، وينبغي أن يقترب مجموعه من
96.
\end{solution}

\begin{solution}{exo:b2:field-lab-methods-2:8}
(أ) $\beta = \ln(0.1194/0.11905)/\ln(85.909261/87.905612) = 0.002936/(-0.02297) = -0.1278$.
(ب) $0.71062 \times (86.908878/85.909261)^{-0.1278} = 0.70957$. (ج) النسبة 86/88 المقيسة أدنى من
النسبة الحقيقية: فقد فُضِّل النظير الثقيل، وجاءت النسبة 87/86 المقيسة أعلى ممّا ينبغي.
\end{solution}

\begin{solution}{exo:b2:field-lab-methods-2:9}
(أ) $0.1000 \times 15.271 \times 10^{-6}/83.9242 \times 1.3599 \times 10^{-3} = \qty{2.474e-11}{mol}$.
(ب) بحسب \eqref{eq:b2:field-lab-methods-2:id}، $N_x = 2.474 \times 10^{-11} \times (732.8 -
0.874)/(0.874 - 0.0568) = \qty{2.216e-8}{mol}$ من $^{86}\mathrm{Sr}$؛ والسترونشيوم: $2.216 \times
10^{-8}/0.0986 \times 87.62 = \qty{1.969e-5}{g}$ في \qty{0.0500}{g}، أي \qty{394}{ppm} (393.9).
(ج) نسبة استرجاع قدرها \qty{99.95}{\percent}: فالطريقة صحيحة ضمن تقريب النسبة.
\end{solution}

\begin{solution}{exo:b2:field-lab-methods-2:10}
(أ) $R_m = \sqrt{732.8 \times 0.05680} = 6.45$، حيث $\mathcal M = 1.018$. (ب) $\mathcal M = 2.31$
عند 0.1، و1.13 عند 0.5، و1.038 عند 1.664 و1.16 عند 100. (ج) قرب $R_x$ لا يكاد المزيج يختلف عن
العيّنة، وقرب $R_s$ عن المُضاف، فيضيع الفرق الذي يحمل المعلومة في
خطأ $R_m$: إذ يتباعد $\mathcal M$ عند الطرفين. وبينهما تكون القيمة الصغرى مسطّحة جدًّا
($\mathcal M < 1.2$ من نحو 0.4 إلى 100)، فلا يحتاج المحلّل إلا إلى رتبة القدر الصحيحة
للمُضاف، وهي ما يوفّره تركيز تقريبي، من فلورة الأشعة السينية مثلًا.
\end{solution}

\begin{solution}{exo:b2:field-lab-methods-2:11}
(أ) $\sigma = \sqrt{0.62 \times 0.38/400} = 0.0243$: أي $62 \pm 2.4$. (ب) $(65 - 62)/2.43 = 1.2$
انحراف معياري: فقد يكون الصخر غرانوديوريتًا فعلًا. (ج) $\sigma = 0.015$ يتطلّب $0.62 \times
0.38/0.015^2 = 1047.1$، أي 1048 نقطة فلسبار.
\end{solution}

\begin{solution}{exo:b2:field-lab-methods-2:12}
(أ) تخطئ كلّ نقطة باحتمال $1 - p$، مستقلّة عن غيرها: $(1-p)^n$. (ب) $0.99^n \le 0.05$
يعطي $n \ge \ln 0.05/\ln 0.99 = 298.1$: أي 299 نقطة. (ج) $n \ge \ln 0.05/\ln 0.9995 = 5990$: فعدّ الشبكة
أداة غير ملائمة؛ ويُبحث عن المعدن الثانوي النادر بمسح الشريحة كلّها.
\end{solution}

\begin{solution}{pb:b2:field-lab-methods-2:1}
\noindent\textbf{الجزء الأول.}
\textbf{1.} الكوارتز $372/1200 = \qty{31.0}{\percent}$، والفلسبار القلوي 29.0، والبلاجيوكلاز 27.0، والبيوتيت
9.0، والمسكوفيت 3.0، والأكاسيد المعتمة 1.0.
\textbf{2.} $\sqrt{0.31 \times 0.69/1200} = 0.0134$؛ والمسكوفيت: $\sqrt{0.97/(0.03 \times 1200)} =
\qty{16.4}{\percent}$.
\textbf{3.} $M = 9 + 3 + 1 = \qty{13.0}{\percent}$؛ و$M' = 13 - 3 = \qty{10.0}{\percent}$ من دون
المسكوفيت: فاتح اللون.
\textbf{4.} $Q + A + P = 1044$ نقطة: \mbox{\textdir TLT Q $= 35.6$}، و\mbox{\textdir TLT A $= 33.3$}، و\mbox{\textdir TLT P $= 31.0$}.
\textbf{5.} $324/672 = \qty{48.2}{\percent}$، مع Q بين 20 و60: إنّه غرانيت، وبين 35
و65 فهو مونزوغرانيت.
\textbf{6.} $\sqrt{0.482 \times 0.518/672} = 0.0193$: يقع الحدّان عند 35 و65 على بعد 6.9 و8.7
انحراف معياري. فالتسمية مؤكّدة.
\textbf{7.} $0.31 \times 0.69/0.01^2 = 2139$ نقطة.

\noindent\textbf{الجزء الثاني.}
\textbf{8.} $0.0051/0.8000 = \qty{0.64}{\percent}$.
\textbf{9.} $0.7949/8.7949 = 0.0904$، و$72 \times 0.0904 = \qty{6.51}{\percent}$ من \ce{SiO2} في
الزجاج: فالصهر يخفّف العيّنة إحدى عشرة مرّة.
\textbf{10.} $\bar I = 24.91$، و$\bar c = 62.283$، و$\sum(I - \bar I)^2 = 31.918$، و$\sum(I - \bar I)(c -
\bar c) = 80.009$: إذن $a = 2.5067$ لكلّ \unit{kcps}، و$b = 62.283 - 2.5067 \times 24.91 = -0.16$.
\textbf{11.} $2.5067 \times 28.63 - 0.16 = \qty{71.61}{\percent}$.
\textbf{12.} BCR-2 \mbox{$-0.04$}، وAGV-2 \mbox{$+0.05$}، وRGM-1 \mbox{$-0.01$}: كلّها أصغر من الانحرافات المعيارية
للطريقة (0.1 إلى 0.2)، فالخطّ ملائم على هذا المدى.
\textbf{13.} $71.61 + 27.15 + 0.64 = \qty{99.40}{\percent}$، بين 97 و101: لا يُعاد.
\textbf{14.} $71.61 \times 100/(71.61 + 27.15) = \qty{72.51}{\percent}$.

\noindent\textbf{الجزء الثالث.}
\textbf{15.} $0.0986/0.8258 = 0.1194$: فقيمة التطبيع هي النسبة الطبيعية.
\textbf{16.} $\beta = 0.007313/(-0.022972) = -0.3183$.
\textbf{17.} $0.71655 \times (86.908878/85.909261)^{-0.3183} = 0.71392$.
\textbf{18.} تبتعد النسبة 86/88 المقيسة بمقدار $0.11853/0.1194 - 1 = \qty{-0.73}{\percent}$ عن القيمة الحقيقية
على وحدتي كتلة: أي نحو \qty{0.36}{\percent} لكلّ وحدة كتلة، لصالح النظائر الثقيلة.
\textbf{19.} نسبة $^{86}\mathrm{Sr}/^{88}\mathrm{Sr}$ في المُضاف 0.7661: فلم تعد النسبة 86/88 في المزيج
القيمة الطبيعية 0.1194، وتثبيتها عند 0.1194 يعامل المُضاف كأنّه انحياز في الجهاز.

\noindent\textbf{الجزء الرابع.}
\textbf{20.} مجموع النسب إلى 88: $561.4 + 0.7661 + 0.1819 + 1 = 563.348$؛ و$561.4/563.348 =
\qty{99.65}{\percent}$ و$0.7661/563.348 = \qty{0.13599}{\percent}$.
\textbf{21.} $0.2500 \times 15.271 \times 10^{-6}/83.9242 = \qty{4.549e-8}{mol}$ من السترونشيوم،
و$\times 1.3599 \times 10^{-3} = \qty{6.186e-11}{mol}$ من $^{86}\mathrm{Sr}$.
\textbf{22.} $561.4/0.7661 = 732.8$ و$0.0056/0.0986 = 0.05680$.
\textbf{23.} $N_x = 6.186 \times 10^{-11} \times (732.8 - 1.6812)/(1.6812 - 0.0568) =
\qty{2.784e-8}{mol}$.
\textbf{24.} $\mathcal M = 1.6812 \times 732.74/(731.12 \times 1.6244) = 1.04$؛ تضيف النسبة
$1.04 \times \qty{0.01}{\percent}$، والمُضاف $0.031/15.271 = \qty{0.20}{\percent}$؛ وبالجمع التربيعي
\qty{0.20}{\percent}، تهيمن عليه معايرة المُضاف.
\textbf{25.} السترونشيوم: $2.784 \times 10^{-8}/0.0986 = \qty{2.824e-7}{mol}$، أي $2.824 \times
10^{-7} \times 87.62 = \qty{2.474e-5}{g}$ في \qty{0.1000}{g}. يحوي غرانيت فاردن
\textbf{\qty[text-series-to-math=true]{247}{ppm}} من السترونشيوم، $\pm 0.5$.
\end{solution}
